\section{Morita rigidity and a realized twisted precursor}
\label{sec:m3-theoretical-retry}

The four open rings retain their general unresolved status. The
following results distinguish an actual auxiliary construction from
necessary module data and from a construction of a source ring.

\subsection{Categorical Morita rigidity of $R8\_3\_928$}
\begin{theorem}
Let $\mathcal C$ be a complex fusion category with based ring
$R8\_3\_928$. Every indecomposable semisimple $\mathcal C$-module
category is equivalent to the regular module category. Consequently
every categorical Morita partner is tensor equivalent to $\mathcal C$.
No rank bound on module categories is assumed.
\end{theorem}
\begin{proof}
In the source basis $g=X_7$ acts by
$(0\ 7)(1\ 2)(3\ 4)(5\ 6)$. The established semionic obstruction
forces the nontrivial pointed $C_2$ associator. This conclusion is
invariant under the spherical/Galois reduction. The module pentagon
therefore forbids a $g$-fixed module simple: such a stabilizer would
trivialize that associator.

For a module simple $m$, write its internal end as $A=\sum_i a_iX_i$.
The complete internal-end domain has the form
$(1,a,a,b,b,c,c,0)$, with $0\le a\le7$, $0\le b\le9$, $0\le c\le11$.
These bounds follow from $a_i\le\lfloor\FPdim(X_i)\rfloor$;
the exact dimension character and strict rational root bounds certify
them. Testing all $960$ vectors by entrywise and rational
positive-semidefinite induction-Gram conditions leaves precisely
\[
 (1,0,0,0,0,0,0,0),\ (1,1,1,1,1,1,1,0),\
 (1,1,1,1,1,2,2,0),\ (1,2,2,2,2,3,3,0).
\]
This is a complete finite domain, not a search over a chosen module rank.

Let $B_{i,n}=[X_i m:n]$. Adjunction gives
$(BB^t)_{ij}=\sum_k a_kN_{jk}^{i}$. Remove the columns of $m,gm$;
the others occur in $g$-pairs. Write $u,d$ for the pair sums and
differences on the three remaining row-pairs. Their Gram matrices are
$P=\sum uu^t$ and $M=\sum dd^t$, where $M=I_3$ and
\[
P\in\left\{I_3,
\begin{pmatrix}9&10&14\\10&13&18\\14&18&25\end{pmatrix},
\begin{pmatrix}13&16&18\\16&21&22\\18&22&25\end{pmatrix},
\begin{pmatrix}13&18&20\\18&25&28\\20&28&33\end{pmatrix}\right\}.
\]
Every matrix has determinant one. Cauchy--Binet says that the squared
integral $3\times3$ minors of a factor of $P$ sum to one. There is a
unique unimodular minor. All minors replacing one of its columns must
vanish, forcing every extra column to be zero. Since $m$ generates the
module category, its columns are nonzero. Exactly three residual pairs
remain, and the module category has rank eight.

The complete integral factorizations, with $u\equiv d\pmod2$ and
$|d_i|\le u_i$, give one factor for each internal end up to column
permutation. The $48$ signed permutation factors of $M=I_3$ exhaust its
possibilities. Each resulting $B$ is invertible, so the source actions
are determined by conjugating the regular matrices. Exact source-product
and dual-transpose checks identify all four with the regular based
module. This finite computation is replayed from the complete bounds.

The categorical conclusion needs one further step. Each such module
contains a simple $n$ with internal end $1$. The internal-end equivalence
then identifies the actual module category with $\mathcal C_1=\mathcal C$.
No uniqueness of module associators is inferred merely from its based
action. Taking dual categories proves the Morita assertion.
\end{proof}

There are exactly four connected separable algebra isomorphism classes
within any hypothetical categorification: they are the internal matrix
algebras of simples, modulo the four right-$g$ orbits. Their four
underlying objects above are distinct. Every dominant tensor functor
from this category is an equivalence onto its fusion image. Indeed its
target is a regular module category; if the target unit corresponds to
$X$, the positive module dimension vector makes the simple corresponding
to $1$ have dimension $1/\FPdim(X)\ge1$. Hence $X$ is invertible, the
unit internal end is $1$, and adjunction gives full faithfulness.
These statements do not supply the hypothetical category.

\subsection{A complex twisted $GH(C_6,2)$ category}
Put $u=6+\sqrt{37}$. The based ring with invertibles $a_j$, $j\in C_6$,
and noninvertibles $a_j\sigma$ has relations
\[
 a_ia_j=a_{i+j},\qquad a_j\sigma=\sigma a_{-j},\qquad
 \sigma^2=1+2\sum_{j\in C_6}a_j\sigma.
\]
Its dimension $6(1+u^2)=444+72\sqrt{37}$ equals that of $R8\_3\_1768$.
The completed rank-eight based-module route forces pointed associator
class $3\in H^3(C_6,\C^\times)$. Its existence is now established,
without a unitarity assertion.

\begin{proposition}
There is a complex fusion category with these rules and pointed
associator class $3$. This does not assert a categorification of
$R8\_3\_1768$.
\end{proposition}
\begin{proof}
Use the explicit Radchenko--Wheeler near-group realization associated
to $\gamma=\left(\begin{smallmatrix}145&12\\12&1\end{smallmatrix}\right)$,
with group $G=C_{12}^2$ and multiplicity $144$; see
\href{https://arxiv.org/abs/2609.21892}{Radchenko--Wheeler, Section 5.3}.
Its bicharacter and Gaussian are
\[
 B((x,y),(a,b))=e^{2\pi i(yb-xa)/12},\qquad
 q(x,y)=(-1)^{x+y}e^{2\pi i(y^2-x^2)/24}.
\]
These are the specified realization's data, not parameters assigned to
an arbitrary near-group category.

Let $H=\{(x,x+6t):x\in C_{12},t\in C_2\}$ and use the cocycle
$\omega((x,t),(a,s))=(-1)^{ta}$. Its commutator is $B|_H$;
the latter has radical $K=\{(2z,2z):z\in C_6\}=H^\perp$.
On the algebraic twisted crossed product of the Leavitt realization,
extend the source automorphisms and noninvertible endomorphism by
\[
 \beta_g(\lambda_h)=B(g,h)^{-1}\lambda_h,
 \qquad \tau(\lambda_h)=q(h)U_h\lambda_h.
\]
On the original algebra their restrictions are $\alpha_g,\rho$.
The original intertwiner identities verify multiplication and covariance:
the cocycle occurs on both sides equally, and $q(h)^2=1$ on $H$.
Coefficient comparison in the crossed-product direct sum gives
$\operatorname{End}(\tau)=\C^\omega[H]$ and six distinct invertible
classes $G/H$. The original complete splitting intertwiners still
decompose $\tau^2$; its duality maps still satisfy both snake identities.
Thus additive idempotent completion gives a rigid semisimple tensor
category with simple unit and finitely many simples. These are the
algebraic identities checked in the accompanying full construction;
no $C^*$ reconstruction or classification of arbitrary spherical
categories is used.

Set $X=\lambda_{(1,1)}$ and $T=\lambda_{(0,6)}$. Then
$X^{12}=T^2=1$ and $TX=-XT$. Hence
$\C^\omega[H]\cong\bigoplus_{\eta\in C_6}M_2(\C)$.
The diagonal-average projector
$p=12^{-1}\sum_lX^l$ is primitive. Its corner
$\sigma=(p,\tau)$ is simple; the six such central blocks are its
invertible translates and $\tau=2\sum_j a_j\sigma$.

For clarity the crucial square multiplicity has an independent tensor
proof. The diagonal subgroup $L=\{(x,x)\}$ is maximal isotropic and
$q|_L=1$. The algebraic isotropic orbifold identities give an actual
$GH(C_{12},1)$ category there. Inclusion of this crossed product into
the twisted $H$ crossed product preserves its complete basic Hom
lists and all splitting intertwiners, hence extends to a tensor functor
on idempotent completions. Its chosen primitive diagonal projector
remains primitive. The old simple square is
$\sigma_0^2=1+\sum_{j\in C_{12}}b_j\sigma_0$.
The functor identifies $b_j$ with $b_{j+6}$ and sends $\sigma_0$ to
$\sigma$. Applying it proves the displayed multiplicity-two rule and
self-duality. The explicit projector-rank calculation gives the same
answer independently.

For $a=\beta_{(1,0)}$, one has
$a^6=\operatorname{Ad}(\lambda_{(6,0)})$ and
$a(\lambda_{(6,0)})=-\lambda_{(6,0)}$. This cyclic obstruction is
class $3$ in $H^3(C_6,\C^\times)$, as required.
\end{proof}

\subsection{The remaining algebra problem}
The old based-module calculation has precisely two rank-eight
completions for an algebra $1+\sigma$. This does not prove that such
an algebra exists in the category just constructed. In particular the
canonical near-group algebra descends to the matrix algebra
$\operatorname{End}(\sigma)$, not to $1+\sigma$.

More explicitly, write $Z=\lambda_{(0,6)}$ and take the normalized
maps
\[
 r=\frac1{\sqrt{12}}\sum_l S_{(l,l)}\lambda_{(l,l)},\quad
 t_0=\frac1{\sqrt{12}}\sum_lT_{(l,l)}p,\quad
 t_1=\frac1{\sqrt{12}}\sum_lT_{(l,l+6)}Zp.
\]
Their primed algebraic duals satisfy $sr=1$, $t_i't_j=\delta_{ij}p$;
$t_0',t_1'$ span $\operatorname{Hom}(\sigma^2,\sigma)$.
After normalizing the unit-channel product to $s$, a multiplication
on $1+\sigma$ is specified by $t=b_0t_0'+b_1t_1'$. Associativity is
exactly
\begin{align*}
 s(t\otimes1)&=s(1\otimes t),\\
 (s\otimes1)+t(t\otimes1)&=(1\otimes s)+t(1\otimes t).
\end{align*}
The first equation is automatic here: cubic rotation is $I_2$.
The second is a quadratic morphism system in two variables, with
coefficients from the explicit finite quantum dilogarithm. Separability
must also be checked. No solution of this remaining system is claimed.
The categorical dimension in the specified Stark realization is
$6-\sqrt{37}$, whereas its FP dimension is $u$; positivity is not
assumed in these algebraic equations.

\subsection{Further route obstructions}
A proper inclusion of fusion categories has FP index at least two.
Indeed each indecomposable component over the subcategory has dimension
at least its dimension: for a simple $X$, its internal end in the
subcategory is a subobject of $XX^*$, and the internal-end module
dimension formula gives that inequality. A proper inclusion has at
least two components. The $SU(4)_8$ exceptional module-category index
over either quartic target is
$24-15\sqrt2+12\sqrt3-9\sqrt6\in(1,2)$, excluding this route.
The indicated $F_4{}_3/G_2{}_4$ routes fail by dimension as well.

For $R8\_3\_1154$, put $d=2+\sqrt5$, $v=d^2$; its dimension is $60v$.
Exact Weyl arithmetic excludes five entire families with arbitrary
pointed factors and one connected étale condensation:
$SU(4)_6$, $SU(5)_5$, $SU(6)_4$,
$SU(3)_7\boxtimes\mathrm{Fib}$, and the rank-fourteen modular factor
constructed in the rank-ten conformal model. In each case a rational
required algebra dimension either violates algebraic integrality or
exceeds the pointed dimension bound: every nonpointed simple has a
strictly positive nonconstant coordinate in the exact positive
dimension basis, so a rational-dimensional algebra must be pointed.
The exact field enumerations and the all-pointed-factor argument are
in the accompanying proof. These are obstructions to specified
construction families, not general noncategorifiability theorems.

All new proofs, complete finite domains and independent audits are
curated in \path{evidence/unresolved_retry_20261024/README.md}.
The ten remaining open cases include these four rank-eight rings.
